% Configuration audit continuation: exact mixed-characteristic frontiers.
% AI-produced by GPT-6 Astra; provisional statement identifiers B.CF12--13.
\subsection{Wild Kummer obstructions and the unbounded nilpotent tower}

Throughout this subsection,
\[
 K_p=\operatorname{Frac}(\Z_{(p)}^{sh})\subset\overline\Q,
 \qquad I=G_{K_p}=I_{\Q_p},\qquad
 F_p=\overline\Q^{I'},\qquad B_p=\overline\Q^{I''}.
                                                        \tag{B.MIX.1}
\]
The identification with local inertia is the henselization--completion
comparison, with the chosen prolongation of the valuation. It concerns
absolute Galois groups; the global field \(K_p\) has not been replaced
by \(\Q_p^{\mathrm{ur}}\). For clarity, a finite collection of global
algebraic numbers and its local Galois conjugates lie in a number field
\(D\) whose completion \(M\) is finite Galois over \(\Q_p\).
Then \(I\cap G_M=I_M\) fixes the compositum \(K_pD\).
These are the identifications used below.

\begin{claim}{B.CF12}{Proved}{Wild covers obstruct supported scaling at the metabelian field}
For a finite extension \(E/B_p\), let \(\kappa_E\) be the rational
\(p\)-primary Kummer map. Put
\[
 W_p=\{a\in B_p^\times:a^m\in F_p^\times
                         \text{ for some }m\geq1\}.
\]
Then
\[
 \ker\kappa_E=\bigcap_{n\geq1}(E^\times)^{p^n}=W_p,
 \qquad W_p\cap(1-W_p)=F_p\setminus\{0,1\}.             \tag{B.MIX.2}
\]
In particular the supported scaling used in B.CF6--7 cannot be the
point function of a pinned relative autoequivalence over \(B_p\).
Precisely, for \(\lambda=1+p\), the bijection
\[
 f(a)=\begin{cases}\lambda a,&v(a)>0,\\a,&v(a)\leq0,
                         \end{cases}\qquad f(0)=0
                                                        \tag{B.MIX.3}
\]
of \(\overline\Q\) is excluded by two one-puncture
\(p\)-Kummer comparisons. This is an obstruction to this proposed
counterexample, and not a proof of reconstruction over \(B_p\).
\end{claim}
\begin{proof}
We independently reprove the kernel input of AS-FF.42. Suppose
\(\kappa_E(a)=0\). Enlarge to \(E'=B_pD\), with \(D\) a number
field stable under the local Galois group, containing \(a\) and the
finitely many generators for \(E/B_p\). Its local completion
\(M/\Q_p\) is finite Galois. Write
\[
 J=I\cap G_M=I_M,\quad H=I''\cap J,
 \quad\Gamma=J/H,\quad C=(I'\cap J)/H.
\]
The Kummer class of \(a\) on \(J\) inflates from
\(H^1(\Gamma,\Q_p(1))\). The group \(C\) and the quotient
\(\Gamma/C\) are abelian; the cyclotomic character is trivial
on \(C\) and has infinite image on \(\Gamma/C\).
For an abelian group \(A\) with nontrivial scalar character \(\chi\),
the cocycle equation gives
\((\chi(g)-1)c(h)=(\chi(h)-1)c(g)\). Choosing one
\(g\) with \(\chi(g)\ne1\) proves \(H^1(A,\Q_p(\chi))=0\).
Inflation--restriction therefore injects
\(H^1(\Gamma,\Q_p(1))\) into \(H^1(C,\Q_p(1))\).
Every \(g\in I'\) centralizes \(C\), because the commutator
lies in \(I''\cap J\), and fixes its coefficient module.
Thus it fixes the class, so
\(\kappa_J(g(a)/a)=0\).

Both \(a\) and \(g(a)\) belong to \(M\). Restriction from
\(H^1(G_M,\Q_p(1))\) to \(H^1(I_M,\Q_p(1))\) is injective,
since \(\Q_p(1)^{I_M}=0\). Over the finite local field \(M\),
the rational Kummer kernel is exactly \(\mu(M)\): valuation
detects the uniformizer component, and the logarithm identifies a
sufficiently small principal-unit group with an additive lattice
in \(M\). Hence \(g(a)/a\in\mu(M)\) for every \(g\in I'\).
An integer \(m\) killing the finite group \(\mu(M)\) makes
\(a^m\) fixed by \(I'\), so \(a^m\in F_p\).
Since \(F_p\) contains all roots of unity,
\(F_p(a)/F_p\) is cyclic and belongs to its maximal abelian
extension \(B_p\). Thus \(a\in W_p\).

Conversely, if \(a^m\in F_p\), every \(p^n\)-th root of
\(a\) is a radical of an element of \(F_p\). Its extension
over \(F_p\) is cyclic, so it belongs to \(B_p\).
Compatible roots exist. Since \(G_E\) fixes \(\mu_{p^\infty}\),
integral Kummer classes are homomorphisms to the torsion-free
module \(\Z_p(1)\); rational vanishing is the same as integral
vanishing. This proves the first equality in (B.MIX.2).
If \(a,1-a\in W_p\), then for each \(g\in I'\) one has
\(g(a)=\zeta a\), \(g(1-a)=\eta(1-a)\), with
\(\zeta,\eta\in\mu_\infty\subset F_p\).
If \(g(a)\ne a\), their sum gives
\(a=(1-\eta)/(\zeta-\eta)\in F_p\), a contradiction;
the case \(\zeta=\eta\) forces both to be one.
This proves the second equality.

Choose \(x\in B_p\setminus F_p\) with \(v(x)>0\), and
\(c\in K_p\) with both \(c\) and \(1-c\) units. Such an
\(x\) exists because \(I'\) is a nontrivial closed subgroup
of the free wild pro-\(p\) group, hence has nontrivial
abelianization; multiplying an element of \(B_p\setminus F_p\)
by a sufficiently large power of \(p\) ensures positive value.
This is also the nontriviality argument in AS-FF.40. The residue
field \(\overline\F_p\) supplies \(c\), including when \(p=2\).
Put \(y=1+c(x-1)\). Then \(y\) is a unit,
\(f(1)=1\), \(f(y)=y\), and \(f(x)=\lambda x\).

Suppose the stated relative autoequivalence existed. On the line
punctured at a point \(a\in B_p\), it sends the rank-one geometric
\(p\)-Kummer line to itself by a scalar
\(\alpha_a\in\Z_p^\times\). The remaining summand is an
arithmetic Kummer class, which cancels on subtracting two point
evaluations. The puncture at \(1\), evaluated at \(y,x\), gives
\[
 \kappa_{B_p}\!\left(\frac{y-1}{\lambda x-1}\right)
     =\alpha_1\kappa_{B_p}\!\left(\frac{y-1}{x-1}\right)
     =\alpha_1\kappa_{B_p}(c)=0.
\]
The puncture at \(x\), evaluated at \(y,1\), similarly gives
\[
 \kappa_{B_p}\!\left(\frac{y-\lambda x}{1-\lambda x}\right)
     =\alpha_x\kappa_{B_p}(1-c)=0.
\]
The two left arguments are \(A\) and \(1-A\), for
\(A=c(x-1)/(\lambda x-1)\). Equation (B.MIX.2) forces
\(A\in F_p\). But this fractional linear transformation has
coefficients in \(F_p\) and nonzero determinant
\(c(\lambda-1)\), so it cannot carry
\(x\notin F_p\) into \(F_p\). This contradiction pinpoints the
wild arithmetic obstruction absent from the residue-zero proof.
\end{proof}


\begin{claim}{B.CF13}{Proved}{A prime-dependent Kummer theorem for the unbounded nilpotent tower}
Let
\[
 N=\bigcap_{c\geq1}\gamma_{c+1}(I),\qquad
 K_{p,\mathrm{nil}}=\overline\Q^N
      =\bigcup_{c\geq1}\overline\Q^{\gamma_{c+1}(I)}.
                                                        \tag{B.MIX.4}
\]
For every \(p\geq5\) and every finite extension
\(E/K_{p,\mathrm{nil}}\),
\[
 \ker\!\left(E^\times\longrightarrow H^1(G_E,\Q_p(1))\right)
       =\bigcap_{n\geq1}(E^\times)^{p^n}
       =\mu_\infty.                                    \tag{B.MIX.5}
\]
For \(p=2\) and \(p=3\), the corresponding kernel is already
strictly larger than \(\mu_\infty\) over
\(K_{p,\mathrm{nil}}\) itself. Explicit nontorsion elements in it are
\[
       2\quad(p=2),\qquad
       \frac{2+\sqrt{-3}}{2-\sqrt{-3}}\quad(p=3).         \tag{B.MIX.6}
\]
These conclusions decide this exact Kummer input to AS-FF.39 at its
unbounded union. They do not decide reconstruction at that union.
\end{claim}
\begin{proof}
Write \(K=K_{p,\mathrm{nil}}\) and \(Q=I/N\).
The quotient \(Q\) is the maximal pronilpotent quotient of
\(I\): a finite quotient is nilpotent exactly when it kills
some lower-central term, and any finite compositum of such
quotients is again nilpotent. Its pro-\(p\) Sylow subgroup
\(Q_{(p)}\) is normal, and
\[
                    Q=Q_{(p)}\times T,
                                                        \tag{B.MIX.7}
\]
where \(T\) is the tame prime-to-\(p\) procyclic group.
Indeed every finite quotient of \(I\) has a normal wild
\(p\)-subgroup and cyclic prime-to-\(p\) quotient, by the
inertia sequence verified in B.CF11. In a nilpotent quotient
the different Sylow subgroups commute. Taking inverse limits
gives (B.MIX.7); the full tame quotient survives, since it
is already abelian. All roots of unity belong to \(K\).
For odd \(p\), the cyclotomic character on \(T\) surjects
onto \(\mu_{p-1}\subset\Z_p^\times\), and its pro-\(p\)
part comes from \(Q_{(p)}\).

Let \(a\in E^\times\) have zero rational Kummer class.
As above, \(G_E\) fixes the coefficient module, so this is
equivalent to having all compatible \(p\)-power roots of
\(a\) in \(E\). Enlarge to a finite extension \(E'/K\)
normal under the local absolute Galois group, containing the
finite local orbit \(a_1=a,\ldots,a_m\). This is possible
by adjoining the finitely many local conjugates of generators
for \(E/K\); the field \(K\) is stable under \(G_{\Q_p}\),
since \(N\) is characteristic in its normal inertia subgroup.
Choose a number field \(D\subset E'\), stable under
\(G_{\Q_p}\), containing these generators and conjugates,
and let \(M=D_v\), a finite Galois extension of \(\Q_p\).
Put
\[
 F=K_pD(\mu_{p^\infty}),\qquad
 R=F(a_j^{1/p^n}:1\leq j\leq m,\ n\geq1).
\]
Then \(R/K_p\) is Galois and \(R\subset E'\).
The group \(D_R=\operatorname{Gal}(R/F)\) embeds by
Kummer characters into \(\Z_p(1)^m\); in particular it is
torsion-free. The image \(C_R\) of \(N\) in
\(\operatorname{Gal}(R/K_p)\) is finite, because
\(E'/K\) is finite.

Choose \(t\in T\) whose cyclotomic character is a primitive
\((p-1)\)-st root \(\zeta\), and choose a lift
\(\tau\in I\) of \((1,t)\in Q_{(p)}\times T\).
The cyclotomic character factors through \(Q\), so
\(\chi(\tau)=\zeta\) exactly. Its image is central modulo
\(C_R\). Consequently
\[
             [\tau,D_R]\subset C_R\cap D_R=1.           \tag{B.MIX.8}
\]
Here the last equality uses finite versus torsion-free, not
an assertion that the entire radical extension is abelian
over \(K_p\).

Set \(J=I_M=I\cap G_M\) and \(J_0=J\cap\ker\chi\).
The group \(J_0\) fixes \(K_p,D,\mu_{p^\infty}\) and
maps onto \(D_R\). Kummer conjugation and (B.MIX.8) give
\[
       \kappa_{J_0}(\tau a)=\zeta\,\kappa_{J_0}(a).
                                                        \tag{B.MIX.9}
\]
There is no lost subgroup here: \(J\) already fixes the
strict-local field \(K_p\). Restriction
\(H^1(J,\Q_p(1))\to H^1(J_0,\Q_p(1))\) is injective,
because \(J/J_0\) is an abelian group with nontrivial
cyclotomic character and its first cohomology vanishes by
the displayed cocycle calculation in B.CF12. Restriction
from \(G_M\) to \(I_M\) is injective as well, since the
inertia invariants of \(\Q_p(1)\) vanish. Both \(a\) and
\(\tau a\) belong to \(M\), so (B.MIX.9) implies
\[
                \kappa_M(\tau a)=\zeta\kappa_M(a).
                                                        \tag{B.MIX.10}
\]

Valuation in (B.MIX.10) gives \(v(a)=\zeta v(a)\);
as \(\zeta\ne1\), the element \(a\) is a unit.
The principal-unit logarithm then gives
\(\log(\tau a)=\zeta\log(a)\).
The opened Baker--Brumer theorem, with proof in
\citep[Théorème, pp.~1--2]{Rousseau1969}, states that
\(p\)-adic logarithms of algebraic units which are linearly
independent over \(\Q\) remain linearly independent over
\(\overline\Q\). If \(\log(a)\ne0\), the preceding
algebraic relation therefore makes the two logarithms
\(\Q\)-dependent, and forces \(\zeta\in\Q\).
For \(p\geq5\), a primitive \((p-1)\)-st root is not
rational. Thus \(\log(a)=0\); the unit \(a\) in the
finite local field \(M\) is a root of unity. Conversely
all roots of unity have all their divisions in \(K\).
This proves (B.MIX.5), with every finite extension included.

For \(p=2\), each splitting field over \(K_2\) of
\(X^{2^n}-2\) has a finite \(2\)-group as Galois group:
adjoin the \(2^n\)-th roots of unity and then a Kummer
root. These are nilpotent groups, so all those roots lie
in \(K_{2,\mathrm{nil}}\), proving the first exception.

For \(p=3\), put \(z=\sqrt{-3}\),
\(F_0=K_3(z)=K_3(\mu_3)\), and
\(a=(2+z)/(2-z)\). The nonidentity element of
\(\operatorname{Gal}(F_0/K_3)\) sends \(a\) to \(a^{-1}\).
Let \(R_n=F_0(\mu_{3^n},b_n)\), where \(b_n^{3^n}=a\).
This is Galois over \(K_3\), and its kernel over \(F_0\)
is a \(3\)-group \(P_n\). Cauchy's theorem supplies an
involution \(s\) in the Galois group. It inverts
\(\mu_{3^n}\). Writing \(s(b_n)=\xi b_n^{-1}\), its
order two forces \(\xi^2=1\), hence \(\xi=1\).
For \(g\in P_n\), write
\(g(b_n)=\eta b_n\) and \(g(\zeta_{3^n})=\zeta_{3^n}^{u}\).
Direct substitution on these two generators gives
\(sgs^{-1}=g\). Thus
\(\operatorname{Gal}(R_n/K_3)=P_n\times\langle s\rangle\)
is nilpotent. All \(3\)-power roots of \(a\) lie in
\(K_{3,\mathrm{nil}}\). Finally
\(a+a^{-1}=2/7\), which is not an algebraic integer,
so \(a\) is not a root of unity. This proves the second
exception and completes the stated Kummer classification.
\end{proof}

\paragraph{Why the new Kummer theorem does not invoke an existing reconstruction theorem.}
The amenability hypothesis in
\citep[Definition~5.15(5.15.1)]{CarlsonHaineWolf2024}
requires each torus to have no nonzero divisible rational
points, with all primes included in characteristic zero.
It fails already for \(\mathbb G_m\), since
\(\mu_\infty\subset K_{p,\mathrm{nil}}\).
The independent condition~(5.15.2), requiring the Tate module
of \(\operatorname{Pic}(X)\) to vanish for every regular
finite-type \(X\), fails as well. For an explicit example take
\(A:y^2+xy=x^3-p\), with \(c_4=1\) and
\(\Delta=p(1-432p)\). Its node has the two rational tangent
lines \(y=0\), \(y+x=0\), hence split multiplicative reduction.
The opened Tate uniformization and torsion sequence
\citep[Theorem~1, author p.~6; pp.~18--19]{Tate1995}
give a compatible torus Tate line \(\Z_p(1)\subset T_pA\).
The subgroup \(I'\) fixes this line pointwise. Torsion points
of the globally defined elliptic curve are algebraic over
\(\Q\), so their coordinates lie in \(F_p\), and consequently
in \(K_{p,\mathrm{nil}}\). Their divisor classes give
\(\Z_p\subset T_p\operatorname{Pic}(A_{K_{p,\mathrm{nil}}})\).
Thus even the Picard condition has an actual counterexample;
it is not merely absent from the proof of (B.MIX.5).

The explicit positive argument for a fixed bounded nilpotency
class in AS-FF.39 also uses an inertia-invariant lift of every
visible generalized-Albanese functional. Passing to a
pronilpotent quotient does not retain its bounded-nilpotency
Lie-algebra argument: arbitrary pro-\(p\) groups are
pronilpotent. Formula (B.MIX.5) alone supplies neither that
lift nor compatible morphisms on all finite étale covers.
The unbounded tower therefore still requires a separate
reconstruction argument. The larger kernels at \(p=2,3\)
likewise refute a particular input, not reconstruction.
No single-topos reconstruction conclusion is deduced from
ordinary pro-reconstruction.

\subsection{Affine point functions over the unbounded tower when \texorpdfstring{\(p\geq5\)}{p >= 5}}

\begin{claim}{B.CF14}{Proved}{An actual Kummer value detects a rational scalar}
Let \(p\geq5\), \(K=K_{p,\mathrm{nil}}\), and let \(E/K\)
be finite. If \(b\in E^\times\) and \(\alpha\in\Q_p\)
satisfy
\[
                 \kappa_E(b)=\alpha\kappa_E(p),           \tag{B.MIX.11}
\]
then \(\alpha\in\Q\).
\end{claim}
\begin{proof}
Pass to a finite extension \(E'/K\) normal under
\(G_{\Q_p}\), as in B.CF13, and put \(N'=G_{E'}\).
It is normal in \(I\), contained in \(N=G_K\), and of
finite index in \(N\). For \(g\in I\), conjugate (B.MIX.11)
on \(N'\). The restricted Kummer class of the rational
number \(p\) is \(I\)-invariant, since it comes from
\(H^1(I,\Q_p(1))\). Hence
\(\kappa_{E'}(g(b)/b)=0\). B.CF13 makes every
\(g(b)/b\) a root of unity. The finite local orbit of
the globally algebraic number \(b\) lies in one number
field, whose roots of unity form a finite group. A single
integer \(m>0\) therefore gives
\(b^m\in\overline\Q^I=K_p\).

The class
\(\kappa_I(b^m)-m\alpha\kappa_I(p)\) restricts to zero
on \(N'\), so it inflates from \(H^1(I/N',\Q_p(1))\).
This group is zero. Indeed the kernel \(N/N'\) of
\(I/N'\to I/N=Q\) is finite and fixes the coefficient
module, so its rational first cohomology is zero.
Inflation--restriction identifies the group with
\(H^1(Q,\Q_p(1))\). The group \(Q\) has a central
tame element \(t\) with \(\chi(t)\ne1\).
For every cocycle \(c\), centrality gives
\((\chi(t)-1)c(g)=(\chi(g)-1)c(t)\), so \(c\) is a
coboundary. This proves the claimed vanishing.

The element \(b^m\in K_p\) lies in a finite unramified
local extension \(M/\Q_p\), with \(I_M=I\).
The local Kummer class comparison is injective upon
restriction to inertia, because \(\Q_p(1)^I=0\).
Thus the preceding zero class already vanishes on
\(G_M\). Normalizing its valuation by \(v_M(p)=1\)
now yields
\(m\alpha=v_M(b^m)\in\Z\). Hence \(\alpha\) is rational.
\end{proof}

\begin{claim}{B.CF15}{Proved}{Affine actual point functions for all pinned affine-line automorphisms}
Let \(p\geq5\), and let \(E\) be any finite extension
of \(K_{p,\mathrm{nil}}\). Every pinned autoequivalence
of \(\mathbb A^1_{E,\mathrm{et}}\) over \(BG_E\), with
its relative structural isomorphism retained, has actual
geometric point function
\[
                         a\longmapsto ua+v,
                 \qquad u\in E^\times,\ v\in E.          \tag{B.MIX.12}
\]
The coefficients are unique. This conclusion concerns the
actual point function; it does not identify the remaining
pointwise-trivial automorphism kernel on the full étale site,
or assert existence for arbitrary noninvertible pinned maps.
\end{claim}
\begin{proof}
Let \(f:\overline\Q\to\overline\Q\) be the actual point
function. Relative point evaluation
\citep[Proposition~3.6]{CarlsonHaineWolf2024}, together with
the inverse equivalence, makes it a \(G_E\)-equivariant
bijection. On every finite extension \(L/E\), it gives a
bijection of \(L\). Compose with the unique affine scheme
automorphism taking \(f(0),f(1)\) to \(0,1\); henceforth
\(f(0)=0\), \(f(1)=1\).

For \(a\in\overline\Q\), the equivalence after a finite
scalar extension restricts to the lines punctured at
\(a\) and \(f(a)\). On their rank-one geometric
\(p\)-Kummer lines it acts by
\(\alpha_a\in\Z_p^\times\). The arithmetic summand
cancels when two point evaluations are subtracted. Thus,
over any common finite extension containing \(a,b,x\),
\[
 \kappa_L\!\left(\frac{f(x)-f(a)}{f(b)-f(a)}\right)
       =\alpha_a\kappa_L\!\left(\frac{x-a}{b-a}\right)
       \qquad (b,x\ne a).                              \tag{B.MIX.13}
\]
This only uses the one-puncture cohomology calculation;
no preservation theorem for the place at infinity and
no equality of different inertia exponents is assumed.
Set \(b=a+1\), \(x=a+p\). B.CF14 applies to the left
argument in (B.MIX.13), and proves
\(\alpha_a\in\Q^\times\) for every \(a\).

Fix a complex embedding of \(\overline\Q\), and write
\(|\cdot|\) for its ordinary complex absolute value.
Clear the rational denominator of \(\alpha_a\) in
(B.MIX.13). B.CF13 says that the resulting multiplicative
error is a root of unity. Its complex absolute value is one.
Consequently there is a positive real constant \(C_a\)
such that
\[
                 |f(x)-f(a)|=C_a|x-a|^{\alpha_a}
                       \qquad(x\ne a).                  \tag{B.MIX.14}
\]
One can take \(C_a=|f(a+1)-f(a)|\).

We spell out how (B.MIX.14) forces the exponents and constants
to coincide. If \(\alpha_a<0\), the sequence \(f(n)\), for
positive rational integers \(n\to\infty\), tends to \(f(a)\).
For another centre \(b\ne a\), its distance from \(f(b)\)
would therefore tend to the finite positive number
\(|f(a)-f(b)|\). The formula at \(b\) makes that distance
tend to either zero or infinity, since \(\alpha_b\ne0\).
This contradiction proves \(\alpha_a>0\) for every \(a\).
The reverse triangle inequality gives
\[
 \left|C_a|n-a|^{\alpha_a}-C_b|n-b|^{\alpha_b}\right|
                      \leq |f(a)-f(b)|.
\]
Divide by the larger power of \(n\) and let \(n\to\infty\).
Unequal positive exponents give a contradiction, so all
\(\alpha_a\) equal one positive number \(\alpha\).
Dividing again by \(n^\alpha\) gives \(C_a=C_b\).
Normalization at zero and one makes the common constant one.
Thus
\[
                       |f(x)-f(y)|=|x-y|^\alpha
                       \qquad(x,y\in\overline\Q).       \tag{B.MIX.15}
\]
The triangle with vertices \(0,1,2\) gives
\(2^\alpha\leq2\), hence \(\alpha\leq1\).
Since \(f\) is bijective, its inverse satisfies the same
identity with exponent \(1/\alpha\), and gives
\(1/\alpha\leq1\). Hence \(\alpha=1\).

An isometry of \(\overline\Q\subset\mathbb C\) fixing
zero and one is the identity or complex conjugation.
Indeed distances to zero and one give \(f(i)=i\) or
\(f(i)=-i\); distances to the three noncollinear points
\(0,1,i\) then determine the image of every point, with
the second case giving the reflected point.

The conjugation case cannot be \(G_E\)-equivariant.
First \(E\) is not separably closed. For \(p\geq5\),
the splitting field of \(X^p-p\) over \(K_p\) has
non-nilpotent Galois group
\(\mathbb F_p\rtimes\mathbb F_p^\times\): the radical
degree is \(p\), the cyclotomic tame degree is \(p-1\),
and the natural scalar action is faithful. Thus \(N\ne1\).
It is a closed subgroup of the free wild pro-\(p\) group,
so is torsion-free and infinite; every finite-index subgroup
\(G_E\) is nontrivial. The field \(E\) carries its
nontrivial rank-one henselian mixed-characteristic valuation.
If complex conjugation centralized \(G_E\), it would
preserve \(E\). The opened
\citep[Lemma~2.1]{KhassaKhanduja2010} makes this automorphism
continuous for the valuation; in rank one it preserves the
valuation ring. Uniqueness of prolongation then makes it
preserve the chosen valuation on \(\overline\Q\), so it
lies in the local decomposition group \(G_{\Q_p}\).
Its element of order two maps trivially to the unramified
quotient \(\widehat\Z\), then trivially to the torsion-free
tame inertia quotient. It would lie in wild pro-\(p\)
inertia, which has no element of order two for \(p\geq5\).
This is impossible. The normalized \(f\) is therefore the
identity. Undoing normalization gives (B.MIX.12).
\end{proof}

\begin{claim}{B.CF16}{Proved}{Actual Kummer values force rational boundary residues}
Let \(p\geq5\), let \(E/K_{p,\mathrm{nil}}\) be finite,
and let \(U\) be a smooth geometrically connected curve over
\(E\), with smooth proper compactification \(\overline U\).
Suppose \(\xi\in H^1(U,\Z_p(1))\) evaluates at every point
\(x\in U(L)\), for every finite extension \(L/E\), to the
Kummer class of an actual element of \(L^\times\).
Every geometric boundary residue of \(\xi\) belongs to
\(\Q\cap\Z_p\).
\end{claim}
\begin{proof}
Split the finitely many boundary points over a finite extension
and fix one of them, \(q\). Riemann--Roch supplies a nonconstant
rational function \(g\) with its only pole at \(q\) and with a
zero at every other boundary point: prescribe those finitely
many zeros and allow a sufficiently large pole at \(q\).
The map \(g:\overline U\to\mathbb P^1\) is finite.
Remove zero, infinity, and all branch values from the target,
obtaining \(V\); its inverse image \(U'\) lies in \(U\)
and is finite étale over \(V\). Corestrict \(\xi|_{U'}\)
to a class \(\eta\in H^1(V,\Z_p(1))\).

Its geometric residues vanish away from zero and infinity.
At a removed branch value the original class has zero
residues at all its preimages. The corestriction residue
is the sum of source residues multiplied by residue-field
degrees, not by ramification indices. This formula follows
at finite Kummer levels from valuation of the field norm,
and then by passage to the compatible inverse limit.
The residue at infinity is precisely the residue at the
sole point \(q\), since \(q\) is rational over the chosen
base. Thus the class has the form
\(\eta=\eta_0+\alpha\kappa(T)\), with \(\eta_0\) arithmetic
and \(\alpha=-\operatorname{res}_q\xi\). This also follows
directly from the Kummer sequence on the punctured rational
line, whose Picard group is zero; zero residues at the
additional punctures remove their basis terms.

At every \(t\in V(E)\), the evaluation of \(\eta\) is
an actual Kummer value. Indeed the fibre is a product of
finite separable extensions, and corestriction evaluates
as the product of the field norms of the actual values
of \(\xi\) on these points. Choose \(t_0\in V(E)\) and
a positive integer \(n\) for which \(p^nt_0\in V(E)\);
only finitely many choices are excluded. Subtracting the
two evaluations cancels \(\eta_0\) and gives an actual
value \(b\in E^\times\) with
\(\kappa_E(b)=n\alpha\kappa_E(p)\).
B.CF14 gives \(n\alpha\in\Q\), hence the desired
rationality of the residue. The argument works for each
boundary point and is unchanged by the finite scalar
extension used to split them.
\end{proof}

\begin{claim}{B.CF17}{Proved}{Polynomial actual point functions for arbitrary affine-line morphisms}
Let \(p\geq5\), and let \(E\) be any finite extension
of \(K_{p,\mathrm{nil}}\). Every pinned geometric morphism
\[
       \mathbb A^1_{E,\mathrm{et}}
              \longrightarrow\mathbb A^1_{E,\mathrm{et}}
\]
over \(BG_E\) has actual geometric point function
\(x\mapsto P(x)\) for a unique polynomial \(P\in E[T]\).
Invertibility of the geometric morphism is not assumed.
This does not identify the whole inverse-image functor
with the functor induced by \(P\), and does not prove
existence for an arbitrary source scheme.
\end{claim}
\begin{proof}
Write \(f:\overline\Q\to\overline\Q\) for the actual
point function supplied by relative pinning. If the
underlying generic point maps to a closed point, continuity
makes the entire image that closed point. Evaluation at
zero makes it an \(E\)-rational point, and the assertion
holds with a constant polynomial. Otherwise the generic
point maps to the generic point. Each geometric fibre
\(S_a=f^{-1}(a)\) is then finite: after a finite constant
extension containing \(a\), it is a proper closed subset
of the source affine line. In particular the image of
\(f\) is infinite.

For \(a\in\overline\Q\), restrict to
\(U_a=\mathbb A^1\setminus S_a\) and pull back the
Kummer class of \(T-a\). Split \(S_a\) after a finite
extension. The resulting class is an arithmetic class
plus
\(\sum_{r\in S_a}q_{a,r}\kappa(T-r)\), because a
localization of the polynomial ring has trivial Picard
group. B.CF16 gives \(q_{a,r}\in\Q\cap\Z_p\).
Subtract the evaluation at any fixed \(x_0\notin S_a\),
and clear these finitely many rational denominators.
B.CF13 identifies the resulting multiplicative error
with a root of unity, over every finite field of
point definition. In a fixed complex embedding, this gives
\[
 |f(x)-a|^2=P_a(x),\qquad
 P_a(z)=C_a^2\prod_{r\in S_a}|z-r|^{2q_{a,r}},
 \quad C_a>0,
 \qquad x\in\overline\Q\setminus S_a.                   \tag{B.MIX.16}
\]
No sign or integrality of the exponents has been assumed.
The function \(P_a\) is positive and real analytic off
\(S_a\), and \(\log P_a\) is harmonic there.

Put \(S=S_0\cup S_1\cup S_i\) and
\(\Omega=\mathbb C\setminus S\). Define a real-analytic
function on this connected open set by
\[
 F(z)=\frac{P_0(z)-P_1(z)+1}{2}
          +i\frac{P_0(z)-P_i(z)+1}{2}.                  \tag{B.MIX.17}
\]
The elementary squared-distance identities show that
\(F(x)=f(x)\) at every algebraic \(x\in\Omega\).
Density and continuity then give
\(|F(z)-a|^2=P_a(z)\) on \(\Omega\setminus S_a\)
for every algebraic \(a\). The function \(F\) is
nonconstant, since otherwise \(f\) would have finite image.

Here is the analytic implication in full. At a point
\(z\in\Omega\), write
\(w=F(z)\), \(A=F_z(z)\), \(B=F_{\bar z}(z)\), and
\(C=F_{z\bar z}(z)\). For all algebraic \(a\ne w\)
for which \(z\notin S_a\), harmonicity yields
\[
 \operatorname{Re}\left(\frac{C}{w-a}
                  -\frac{AB}{(w-a)^2}\right)=0.         \tag{B.MIX.18}
\]
There is at most one excluded algebraic centre: all
\(S_a\) consist of algebraic points and are disjoint.
Thus density extends (B.MIX.18) to every complex
\(a\ne w\). Substituting \(b=1/(w-a)\) gives
\(\operatorname{Re}(Cb-ABb^2)=0\) for every nonzero
\(b\in\mathbb C\). The two polynomial coefficients
must vanish, so \(C=0\) and \(AB=0\).
Consequently \(F\) is harmonic and conformal.
More explicitly, if \(F_z\) is nonzero anywhere,
\(F_{\bar z}\) vanishes on an open neighbourhood, and
real-analytic continuation makes it vanish on all of
\(\Omega\). Otherwise \(F_z\) vanishes identically.
Thus \(F\) is holomorphic or antiholomorphic on \(\Omega\).

For \(r\in S\), choose an algebraic centre
\(a\ne f(r)\). Its fibre omits \(r\), so \(P_a\)
is finite and positive near \(r\). Equation (B.MIX.16)
makes \(F\) bounded on a punctured neighbourhood.
The singularity is removable. Choose three noncollinear
such centres; their distance identities at \(r\),
including the original identities for \(f(r)\), show
that the extended value \(F(r)\) equals \(f(r)\).
This gives an entire holomorphic or antiholomorphic
function agreeing with \(f\) on every algebraic point.
For any fixed centre \(a\), (B.MIX.16) bounds its growth
at infinity by a finite power of \(|z|\). Cauchy's
estimates make it a polynomial or an antipolynomial:
\(F(z)=P(z)\) or \(F(z)=P(\bar z)\).
Interpolation at distinct rational integers, whose
images lie in \(E\), gives \(P\in E[T]\) in either case.

The nonconstant antipolynomial case is impossible.
Let \(c\) denote complex conjugation on \(\overline\Q\).
For \(g\in G_E\), equivariance would give
\(P(cg x)=P(gc x)\) for every \(x\). Thus the field
commutator \(h=cgcg^{-1}\) satisfies \(P(hz)=P(z)\)
for every \(z\). Apply this to \(z+n\), with
\(n\in\Z\). If \(d>0\) is the degree of \(P\),
comparison of the coefficients of \(n^{d-1}\) in
\(P(hz+n)=P(z+n)\) gives \(hz=z\).
So \(c\) centralizes \(G_E\), which is excluded by
the valuation and ramification argument in B.CF15.
Hence \(F=P\), proving the asserted actual point
function. Its polynomial is unique because the infinite
set of rational inputs already determines it.
\end{proof}
